\section*{第 \ref{ch:b3:hydrogen-atom} 章——氢原子}

\begin{solution}{exo:b3:hydrogen-atom:1}
(a) $u'' = (1/a^2 - 2/ar)u$；$1/r$ 项匹配的条件为
$\hbar^2/m_{\text{e}}a = k$，即 $a = a_0$，常数项
给出 $E = -\hbar^2/2m_{\text{e}}a_0^2 = -E_{\text{I}}$。(b) $C =
2/a_0^{3/2}$。(c) $[\hbar^2/mk] = \unit{m}$；$[mk^2/\hbar^2] =
\unit{J}$。(d) 由
\cref{exo:b3:schrodinger-three-dimensions:12}的变分论证表明，$-E_{\text{I}}$
是任一归一化态所能达到的最小能量：
不确定性原理为原子划定了能量下界。
\end{solution}

\begin{solution}{exo:b3:hydrogen-atom:2}
(a) \qty{121.6}{nm}；极限为 \qty{91.2}{nm}。(b) $656.3$（红色），
$486.1$（蓝绿色），\qty{434.0}{nm}（紫色）；极限
为 \qty{364.6}{nm}——前三条可见。(c)
自 \qty{1875}{nm} 降至 \qty{820}{nm}，属近红外。(d)
$R_{\text{H}}c\,(1/109^2 - 1/110^2) = \qty{5.01}{GHz}$。
\end{solution}

\begin{solution}{exo:b3:hydrogen-atom:3}
(a) $\dd(r^2\eu^{-2r/a_0})/\dd r = 0$ 在 $r = a_0$ 处成立。(b) $\langle r\rangle = (4/a_0^3)\int_0^\infty
r^3\eu^{-2r/a_0}\dd r = (4/a_0^3)\times 3!\,(a_0/2)^4 = \tfrac32
a_0$。(c) 取 $x = 4$：
$\tfrac12(16 + 8 + 2)\eu^{-4} = 0.24$：约四分之一的时间里，
电子位于 $2a_0$ 之外。(d) 电子并无确定半径：$P(r)$
是一个具有众数、均值与长尾的分布——原子
在因子二的量级上是模糊的。
\end{solution}

\begin{solution}{exo:b3:hydrogen-atom:4}
(a) $l = 0$：一个；$l = 1$：三个；$l = 2$：五个——共九个。(b) $2$，
$8$，$18$。(c) 恰为前三个周期的长度：
元素周期表正是这些壳层的填充记录。(d)
$m$ 简并依然保留（空间仍保持各向同性）；
$l$ 简并被打破，因为价电子所感受的屏蔽势
已不再是纯 $1/r$ 势——
见\cref{exo:b3:hydrogen-atom:8}。
\end{solution}

\begin{solution}{exo:b3:hydrogen-atom:5}
(a) 类氢离子，取 $Z - 1$：$h\nu = (Z-1)^2E_{\text{I}}(1 -
\tfrac14)$。(b) $784 \times \qty{10.2}{eV} = \qty{8.0}{keV}$——
与所测铜的 K$\alpha$ 线相符。(c) 元素排序所依据的整数
是\emph{核电荷}，而非原子量：
莫塞莱谱线中的空隙定位出了尚未发现的元素。(d)
$42^2 \times 10.2 = \qty{18.0}{keV}$——锝，
数十年后被发现，为这一排序提供了确证。
\end{solution}

\begin{solution}{exo:b3:hydrogen-atom:6}
(a) 两个等质量粒子：$\mu = m_{\text{e}}/2$；结合能
\qty{6.8}{eV}，半径 $2a_0 = \qty{0.106}{nm}$。(b) $\tfrac34
\times 6.8 = \qty{5.1}{eV}$：\qty{243}{nm}。(c) 两个光子，能量均为
\qty{511}{keV}，背对背射出（静止时动量守恒）。(d)
正电子发射断层成像：共线的一对
\qty{511}{keV} 光子，正是每个 PET 环用于三角定位的信号。
\end{solution}

\begin{solution}{exo:b3:hydrogen-atom:7}
(a) $a_{\text{ex}} = 0.0529 \times 12.9/0.058 = \qty{11.8}{nm}$。
(b) $E = 13.6 \times 0.058/12.9^2 = \qty{4.7}{meV}$。(c)
自然激子尺寸超过量子点：决定态的是势壁，而非吸引作用——
$1/R^2$ 胜过 $1/R$，正如所预言的。(d)
$k_{\text{B}}T$ 在室温下约为结合能的五倍：
激子会电离；激子物理仅在 $\sim\qty{50}{K}$ 以下
的洁净晶体中出现。
\end{solution}

\begin{solution}{exo:b3:hydrogen-atom:8}
(a) $3s$：其最内层瓣伸入原子实内部，此处
核电荷 $+11$ 几乎未被屏蔽。(b) $-13.6/9 = \qty{-1.51}{eV}$
与 $3d$ 几乎完全相符：$3d$ 轨道位于原子实之外，
感受到净电荷 $+1$；$3s$（及部分 $3p$）感受较深的势而
下沉。(c) $\Delta E = \qty{2.10}{eV}$：$\lambda = \qty{590}{nm}$
——路灯的钠黄光。(d) $n_{\text{eff}} =
\sqrt{13.6/|E|}$：$1.63$ 与 $2.11$，因此 $\delta_s = 1.37$，
$\delta_p = 0.89$。
\end{solution}

\begin{solution}{exo:b3:hydrogen-atom:9}
(a) $n = 100$ 时的结合能约为 $\sim\qty{1.4}{meV}$，再加上
电子较小的热能：一个远红外或无线电光子。(b)
$\Delta n = 1$ 的跃迁（在 $n \approx 100$ 附近）落在 GHz 无线电
波段，即复合谱线。(c) H$\alpha$ 是 $3 \to
2$ 的跃迁，位于级联末端附近。(d) 莱曼 $\alpha$ 线虽最强，
却属紫外线，且在充满基态氢的星云中
会被共振散射而囚禁；H$\alpha$ 则可自由逸出，
使星云呈现红色。
\end{solution}

\begin{solution}{exo:b3:hydrogen-atom:10}
(a) $\langle 1/r\rangle = 1/a_0$：$\langle E_p\rangle = -k/a_0 =
\qty{-27.2}{eV} = 2E_1$。(b) $\langle E_k\rangle = E_1 - \langle
E_p\rangle = +\qty{13.6}{eV}$：维里比 $-\tfrac12$ 适用于
每一条 $1/r$ 轨道。(c) $v = \sqrt{2E_k/m_{\text{e}}} = \alpha c =
\qty{2.2e6}{m/s}$。(d) $I = ev/2\pi a_0 \approx \qty{1}{mA}$
以 \qty{53}{pm} 的半径环流：$B \sim \mu_0I/2a_0 \approx
\qty{12}{T}$——这一巨大的内部磁场，超精细结构
将依赖于它。
\end{solution}

\begin{solution}{exo:b3:hydrogen-atom:11}
(a) 由\cref{exo:b3:hydrogen-atom:10}(c)可得 $v/c = \alpha$；
标度为 $Z\alpha/n$。(b) $\alpha^2E_{\text{I}} = \qty{0.72}{meV}$；
当 $n = 2$ 时，分裂量级为数十 $\unit{\micro eV}$——
约 \qty{10}{GHz}，可用无线电测量（且已有实测）。(c) $Z\alpha =
0.67$：约为光速的三分之二——薛定谔的处理
失效，狄拉克方程取而代之。(d) 当 $Z\alpha \to 1$ 时
基态能量在形式上会跌破 $-2m_{\text{e}}c^2$：
真空本身将激发电子–正电子对的产生——超临界
场，正是在重离子碰撞中被探寻的对象。
\end{solution}

\begin{solution}{exo:b3:hydrogen-atom:12}
(a) 求导：$\dot{\vect p}\wedge\vect L = -(mk/r^3)\vect
r\wedge(\vect r\wedge\dot{\vect r}\,m)$；展开二重
叉积恰得 $mk\,\dd\vect e_r/\dd t$，于是 $\dot{\vect A} =
\vect 0$。(b) 守恒轴是各向同性之外的又一对称性；
简并是标记不完备，而额外的标记（$\vect
A$ 的量子对应量）使各 $l$ 在同一 $n$ 之内相互连接。(c) 若
$V = -k/r + \epsilon/r^2$，则有效角动量发生偏移，
角向周期不再与径向周期一致，且
椭圆的轴以速率 $\propto\epsilon$ 转动。(d) 相对论
修正在氢原子自身中扮演 $\epsilon$ 的角色：
\cref{exo:b3:hydrogen-atom:11} 中的精细结构，恰是
$l$ 简并的破缺。
\end{solution}

\begin{solution}{pb:b3:hydrogen-atom:1}
\textbf{1.} $\langle r\rangle \sim n^2a_0$；$|E_n| =
E_{\text{I}}/n^2$；能级间距 $\approx 2E_{\text{I}}/n^3$；轨道
频率 $\propto 1/n^3$（玻尔轨道上的开普勒定律）。
\textbf{2.} $E_{n+1} - E_n \to 2E_{\text{I}}/n^3$，且 $h$ 乘以
经典频率 $v_n/2\pi r_n$ 的积等于同一表达式
（\cref{pb:b3:hamiltonian-mechanics:1} 第 20 项的计算）。
\textbf{3.} $n = 50$：$\qty{132}{nm}$，\qty{5.4}{meV}；$n = 100$：
$\qty{0.53}{\micro m}$，\qty{1.4}{meV}。
\textbf{4.} $n^2ea_0 = \qty{2.1e-26}{C.m} \approx 6300\,\text{D}$：
相当于三千个水分子的偶极矩，全部集中在一个原子上。
\textbf{5.} $F \sim \num{5.4e-3}/\num{1.32e-7} \approx
\qty{4e4}{V/m} = \qty{400}{V/cm}$；精确的经典电离
阈值约小八倍（$\sim\qty{50}{V/cm}$）：
无论如何，都只是一缕微弱的场。
\textbf{6.} 取几何截面 $\sim\pi n^4a_0^2$，即便
超高真空的实验室也会使此类原子在
数微秒内发生碰撞；星际密度（$\sim\num{e6}$ 个粒子每
$\unit{m^3}$）时，则使其存活数天——太空才是更好的真空室。
\textbf{7.} $\nu = R_{\text{H}}c\,(1/109^2 - 1/110^2) =
\qty{5.01}{GHz}$。
\textbf{8.} 厘米波射电：一台射电望远镜——大口径天线与
低噪声接收机。
\textbf{9.} $\lambda = \qty{6.0}{cm}$；$n = 110$ 的原子
约为 $\sim\qty{0.6}{\micro m}$：每个波长内
约含十万个原子——完全属于“天线”区间。
\textbf{10.} $\Delta\nu/\nu \sim v/c = \num{4.3e-5}$：约
\qty{200}{kHz}（在 \qty{5}{GHz} 上）——而实测线宽
反过来给出星云的温度。
\textbf{11.} 银河系内侧的电离区：HII 区隐藏在
吸收每个巴耳末光子的尘埃之后——射电复合线
描绘出光学天文学无法看到的旋臂结构。
\textbf{12.} 邻近离子的微电场（斯塔克
展宽）使能级模糊，黑体辐射与宇宙辐射连同
碰撞使这些脆弱的态电离或发生 $l$ 混合。
\textbf{13.} $C_6 \sim d^4/\Delta E \propto (n^2)^4/n^{-3} =
n^{11}$：将 $n$ 由 $1$ 提高到 $70$，相互作用
增长了二十个数量级。
\textbf{14.} 在 $R_{\text{b}}$ 之内，原子对态偏离
共振，因此一个原子的激发\emph{制约}其
邻原子的响应：恰为双量子比特
门所需的受控逻辑。
\textbf{15.} 微米量级的相互作用使每个原子各居于
自己可寻址的光镊中，按原子尺度看相距甚远，
却强耦合——相互作用范围与光学
分辨率相匹配。
\textbf{16.} 每个寿命内约 $\sim 200$ 次门操作：该比值
决定可达到的保真度，起决定作用的是该比值，而非
原始微秒数，才是处理器赖以运转的依据。
\textbf{17.} 在 \qty{300}{K} 下，光子浴
在 $\sim\qty{0.1}{meV}$ 处很稠密：它在
里德伯能级的辐射寿命之内打乱并电离；低温
（\qty{4}{K}）屏蔽则令这些跃迁无从发生。
\textbf{18.} 跃迁速率按 $d^2 \propto n^4$ 标度：当 $n =
50$ 时，约为普通原子耦合的六百万倍——这正是
里德伯原子蒸气池如今能成为经标定的微波
与太赫兹场传感器的原因。
\textbf{19.} $92^2 \times \qty{13.6}{eV} = \qty{115}{keV}$，降至
$13.6/10^6\,\unit{eV} = \qty{13.6}{\micro eV}$：跨越十个数量级，
全由同一公式得出。
\textbf{20.} 能级间距趋于能量的一个
可忽略比例，即 $n \gg 1$ 时，经典性开始显现——
这是原子自身的对应极限。
\textbf{21.} 它是唯一能从第一性原理算到
实验精度的原子：任何偏差都意味着发现，
因此氢原子锚定了基本常数。
\textbf{22.} $\tau \approx 1.6\,\unit{ns} \times 50^3 =
\qty{0.2}{ms}$——颇为充裕，但黑体辐射的重分布（
第三部分）会侵蚀它，这又是需要低温屏蔽的另一理由。
\textbf{23.} CPT 对称性——物质与反物质原子
必须逐条谱线相符；氢原子之所以成为锚点，是因为只有在这里
理论才能与实验一道精确到第十五位数字。
\textbf{24.} 偶极矩正比于 $n^2$，使原子感受到单个
微波光子的场，而其长寿命与能级
结构使它能获得可测相位，\emph{而不}
吸收光子：量子非破坏性传感。
\textbf{25.} 标度 $n^2$、$n^{-2}$、$n^{-3}$、$n^{11}$ 把
这唯一已解的原子从 \qty{52.9}{pm} 拉伸到半微米，
又从 \qty{121.6}{nm} 拉伸到 \qty{6}{cm}——同一个薛定谔
方程之解，既从银河系的 HII 区广播出来，又在光镊阵列中
驱动双量子比特门。
\end{solution}
